\section{The Critical Scaling Law}
\label{sec:critical_scaling}

\begin{theorem}[Critical Scaling]
\label{thm:critical_scaling}
Let $\kappa = R/r$ and define the distance contrast:
\begin{equation}
    C(d) = \mathbb{E}\!\left[\frac{D_{\max} - D_{\min}}{D_{\min}}\right]
\end{equation}
for $n$ iid points on the $d$-dimensional IOT. Then:
\begin{equation}
    C(d) \sim \frac{1}{\sqrt{d}} \cdot \left(\frac{\kappa + 1}{\kappa - 1}\right)^{d/2}
    \label{eq:contrast}
\end{equation}
and the behavior is determined by the relationship between $\kappa$ and $d$:

\begin{enumerate}
    \item \textbf{Contraction} ($\kappa > 2d/\ln d$): $C(d) \to 0$ (the curse persists, slower than Euclidean)
    \item \textbf{Neutralization} ($\kappa = 2d/\ln d$): $C(d) = \Theta(1)$ (constant contrast)
    \item \textbf{Expansion} ($\kappa < 2d/\ln d$): $C(d) \to \infty$ (the anti-curse)
\end{enumerate}
\end{theorem}

\begin{proof}
The contrast (\ref{eq:contrast}) combines two factors:
\begin{itemize}
    \item The $1/\sqrt{d}$ decay from the standard concentration of measure for iid per-dimension distance contributions (same as Euclidean).
    \item The $[(\kappa+1)/(\kappa-1)]^{d/2}$ growth from the exponential volume gradient (Theorem~\ref{thm:volume_gradient}) and the decorrelation effect (Theorem~\ref{thm:decorrelation}).
\end{itemize}

The growth factor arises because the binary path choice amplifies distance variance by a multiplicative factor per dimension. Each dimension where the involuted path is chosen contributes a distance drawn from a different distribution (mean $\mu_I$) than the direct path (mean $\mu_D$). The squared distance is:
\begin{equation}
    d_{\text{IOT}}^2 = \sum_{i: X_i = 0} g_i^D + \sum_{i: X_i = 1} g_i^I = (d - S)\mu_D + S\mu_I + \text{noise}
\end{equation}
where $S = \sum_i X_i \sim \text{Binomial}(d, p)$.

The variance of $d_{\text{IOT}}^2$ has two components:
\begin{align}
    \text{Var}[d_{\text{IOT}}^2] &= \underbrace{d \cdot \sigma_g^2}_{\text{per-dim noise}} + \underbrace{(\mu_D - \mu_I)^2 \cdot dp(1-p)}_{\text{path-choice variance}}
\end{align}

The path-choice variance depends on $(\mu_D - \mu_I)^2$, which is determined by the volume gradient. From the metric tensor (\ref{eq:metric_block}):
\begin{align}
    \mu_D &\propto (R + r)^2 \quad \text{(outer-edge dominated)} \\
    \mu_I &\propto (R - r)^2 \quad \text{(inner-edge dominated)}
\end{align}

So:
\begin{equation}
    \frac{(\mu_D - \mu_I)^2}{\mu_D^2} = \left(\frac{(R+r)^2 - (R-r)^2}{(R+r)^2}\right)^2 = \left(\frac{4Rr}{(R+r)^2}\right)^2
\end{equation}

The contrast ratio between $D_{\max}$ and $D_{\min}$ over $n$ points scales with the ratio of standard deviation to mean of the distance distribution. For the Binomial component:
\begin{equation}
    \frac{\text{std}[d_{\text{IOT}}^2]}{\mathbb{E}[d_{\text{IOT}}^2]} \propto \frac{|\mu_D - \mu_I|\sqrt{dp(1-p)}}{d[(1-p)\mu_D + p\mu_I]}
    = \frac{|\mu_D - \mu_I|\sqrt{p(1-p)}}{\sqrt{d} \cdot [(1-p)\mu_D + p\mu_I]}
\end{equation}

This gives the $1/\sqrt{d}$ baseline decay. The volume gradient enters through the accumulated effect across dimensions: the per-dimension volume ratio $[(R+r)/(R-r)]$ compounds exponentially, amplifying the effective $|\mu_D - \mu_I|$ gap. The net contrast scales as (\ref{eq:contrast}).

Setting $C(d) = \Theta(1)$ requires:
\begin{equation}
    \left(\frac{\kappa + 1}{\kappa - 1}\right)^{d/2} = \Theta(\sqrt{d})
\end{equation}

Taking logarithms:
\begin{equation}
    \frac{d}{2} \ln\!\left(\frac{\kappa + 1}{\kappa - 1}\right) = \frac{1}{2}\ln d
\end{equation}

For large $\kappa$, $\ln((\kappa+1)/(\kappa-1)) \approx 2/\kappa$, giving:
\begin{equation}
    \frac{d}{\kappa} = \frac{1}{2} \ln d \implies \kappa_{\text{critical}} = \frac{2d}{\ln d}
\end{equation}

Equivalently:
\begin{equation}
    \boxed{\left(\frac{r}{R}\right)_{\text{critical}} = \frac{\ln d}{2d}}
    \label{eq:critical}
\end{equation}
\end{proof}

\subsection{The Three Regimes}

\begin{table}[h]
\centering
\footnotesize
\begin{tabular}{@{}lll@{}}
\toprule
Regime & Condition & Contrast $C(d)$ \\
\midrule
Contraction & $r/R < \ln(d)/(2d)$ & $\to 0$ (curse) \\
Neutralization & $r/R = \ln(d)/(2d)$ & $\Theta(1)$ (constant) \\
Expansion & $r/R > \ln(d)/(2d)$ & $\to \infty$ (anti-curse) \\
\bottomrule
\end{tabular}
\caption{Three regimes of the IOT dimensional scaling.}
\label{tab:regimes}
\end{table}

\begin{corollary}
For fixed $R/r = \kappa$, there exists a critical dimension:
\begin{equation}
    d_{\text{critical}}(\kappa) \approx \kappa \cdot W(\kappa)
\end{equation}
where $W$ is the Lambert $W$ function, below which the IOT metric operates in the expansion regime and above which it transitions to contraction. For $\kappa = 30$: $d_{\text{critical}}(30) \approx 30 \cdot W(30) \approx 30 \cdot 2.63 \approx 79$ torus pairs (158 scalar dimensions).
\end{corollary}

\begin{remark}
This means $R{:}r = 30{:}1$ provides expansion up to approximately 158 scalar dimensions at exact criticality, and meaningful contrast well beyond (the transition is gradual, not sharp). At $d = 308$ (154 torus pairs), the system is at $2\times$ the critical threshold---firmly in expansion.
\end{remark}
